% ============================================================
% SECTION 7: MEDIEVAL EUROPEAN MATHEMATICS (~1200 - 1500)
% Slides 51-55
% ============================================================

% ---- Section Divider ----
{
\setbeamercolor{background canvas}{bg=medievalPurple}
\begin{frame}[plain]
  \centering
  \vfill
  \textcolor{white}{\Huge\bfseries Europe Awakens}\\[0.7em]
  \textcolor{white!85}{\Large Arabic Mathematics Crosses the Mediterranean}\\[1.5em]
  \textcolor{white!65}{\large 1200--1500 CE}
  \vfill
  \visualplaceholder[5cm]{Diagram}
  \vfill
\end{frame}
}

\section{Medieval European Mathematics}

% ---- Slide 52: Fibonacci ----
\begin{frame}{Fibonacci --- Bringing Arabic Mathematics to Europe}
  \begin{columns}[T]
    \column{0.55\textwidth}
    \begin{personbox}[Leonardo of Pisa (Fibonacci)]{medievalPurple}
      \textbf{Full name:} Leonardo Bonacci (``Fibonacci'')\\[3pt]
      \textbf{Dates:} c.\ 1170 -- c.\ 1240--50 CE\\
      \textbf{Origin:} Pisa, Italy; traveled extensively in North Africa
    \end{personbox}
    \vspace{0.3em}
    \begin{itemize}
      \item Wrote \textit{Liber Abaci} (1202) --- introduced the \textbf{Hindu-Arabic numeral system} to Europe
      \item The ``Fibonacci sequence'' appeared as a \textbf{rabbit population problem}
      \item Convinced European merchants to abandon Roman numerals
    \end{itemize}

    \vspace{0.3em}
    \begin{insightbox}
      {\footnotesize He did \emph{not} invent the Fibonacci sequence or Hindu-Arabic numerals --- he was the \textbf{messenger} who brought them to Europe.}
    \end{insightbox}

    \column{0.42\textwidth}
    % Fibonacci sequence and nature
    \visualplaceholder[5cm]{Sequence display}

    \vspace{0.3em}
    \visualplaceholder[3.5cm]{Page from \textit{Liber Abaci}\\showing Hindu-Arabic numerals}
  \end{columns}

  \note{
    \textbf{Talking point:} ``He did not invent the Fibonacci sequence or the Hindu-Arabic
    numerals --- he was the messenger who brought them to Europe.''

    The Fibonacci sequence was known to Indian mathematicians (Virahanka, Hemachandra)
    well before Fibonacci.

    \verifytag{Death date very uncertain --- some sources say c.\ 1240, others c.\ 1250.}

    \verifytag{Verify standard attributions: the sequence was known to Virahanka (c.\ 700 CE)
    and Hemachandra (c.\ 1150 CE) before Fibonacci.}
  }
\end{frame}

% ---- Slide 53: Nicole Oresme ----
\begin{frame}{Nicole Oresme --- Graphs Before Descartes}
  \begin{columns}[T]
    \column{0.52\textwidth}
    \begin{personbox}[Nicole Oresme]{medievalPurple}
      \textbf{Dates:} c.\ 1320 -- 1382 CE\\
      \textbf{Origin:} Normandy, France; worked in Paris
    \end{personbox}
    \vspace{0.3em}
    \begin{itemize}
      \item Pioneered \textbf{graphical representations} of changing quantities (``latitude of forms'')
      \item A precursor to \textbf{Cartesian coordinates} --- 250 years before Descartes
      \item Proved the \textbf{harmonic series diverges}:
        \[ 1 + \frac{1}{2} + \frac{1}{3} + \frac{1}{4} + \cdots = \infty \]
      \item Discussed \textbf{fractional exponents}
    \end{itemize}

    \column{0.45\textwidth}
    % Oresme's latitude of forms vs. modern graph
    \visualplaceholder[5cm]{Oresme's version}

    \vspace{0.5em}
    % Comparison arrow
    \visualplaceholder[5cm]{Diagram}

    \vspace{0.2em}
    % Modern coordinate version (tiny)
    \visualplaceholder[5cm]{Diagram}
  \end{columns}

  \note{
    Oresme's ``latitude of forms'' is a remarkable anticipation of coordinate graphing.
    He represented varying quantities (like velocity over time) as areas under curves ---
    the same idea that would later become integral calculus.

    \verifytag{Whether Oresme's proof of harmonic series divergence is the \emph{first} is
    debated. Some attribute priority to others. Say ``one of the first'' to be safe.}
  }
\end{frame}

% ---- Slide 54: Pascal ----
\begin{frame}{Blaise Pascal --- Probability and the Triangle}
  \begin{columns}[T]
    \column{0.50\textwidth}
    \begin{personbox}[Blaise Pascal]{medievalPurple}
      \textbf{Dates:} 1623 -- 1662 CE\\
      \textbf{Origin:} Clermont-Ferrand, France
    \end{personbox}
    \vspace{0.3em}
    \begin{itemize}
      \item Co-founded \textbf{probability theory} (correspondence with Fermat, 1654)
      \item ``Pascal's triangle'' --- known centuries earlier in China, India, and the Islamic world
      \item Built one of the first \textbf{mechanical calculators} (the Pascaline, 1642)
      \item Major philosopher and writer (\textit{Pens\'{e}es})
    \end{itemize}

    \vspace{0.3em}
    \begin{insightbox}
      {\footnotesize ``Pascal's triangle'' was known in China 400 years before Pascal was born. But Pascal discovered deep new properties --- including its connection to probability.}
    \end{insightbox}

    \column{0.47\textwidth}
    % Pascal's triangle
    \visualplaceholder[5cm]{Triangle rows}

    \vspace{0.3em}
    \visualplaceholder[3.8cm]{The Pascaline calculator\\(Mus\'{e}e des Arts et M\'{e}tiers, Paris)}
  \end{columns}

  \note{
    \textbf{Talking point:} ``Pascal's triangle? It was known in China 400 years before
    Pascal was born. But Pascal discovered deep new properties --- including its
    connection to probability.''

    Pascal bridges the Medieval and Early Modern periods. He is placed here for narrative
    flow: connecting to the earlier Yang Hui triangle and al-Karaji's binomial work,
    while setting up the probability work with Fermat.
  }
\end{frame}

% ---- Slide 55: The Stage is Set ----
\begin{frame}{The Stage is Set}
  \begin{columns}[T]
    \column{0.45\textwidth}
    {\small By 1500, the ingredients for a mathematical revolution were in place:}
    \vspace{0.4em}
    \begin{enumerate}\setlength\itemsep{4pt}
      \item[\textcolor{medievalPurple}{\textbf{1.}}] {\small Indian and Islamic mathematics absorbed via Arabic translations}
      \item[\textcolor{medievalPurple}{\textbf{2.}}] {\small The \textbf{printing press} spreading knowledge (Gutenberg, c.\ 1440)}
      \item[\textcolor{medievalPurple}{\textbf{3.}}] {\small Universities growing across Europe}
      \item[\textcolor{medievalPurple}{\textbf{4.}}] {\small A new confidence in human reason (the Renaissance spirit)}
    \end{enumerate}

    \vspace{0.5em}
    {\small\itshape ``And then came the Renaissance, when mathematics exploded\ldots''}

    \column{0.52\textwidth}
    % Timeline of key translations and publications
    \visualplaceholder[5cm]{Vertical timeline}
  \end{columns}

  \note{
    This is a transition slide. The key message: mathematics did not develop in a vacuum.
    It required the confluence of multiple traditions (Indian, Islamic, Greek via Arabic)
    and enabling technologies (the printing press).

    The next section covers the explosive growth from 1500--1700.
  }
\end{frame}
